\documentclass[11pt]{scrartcl}

%%%%%%%%%%%%%% LATEX SAMPLE FILE %%%%%%%%%%%%%%%%
% A line which starts with a % sign
% is called a COMMENT. It is IGNORED
% by the LaTeX processor.

% Include math
\usepackage{amsmath,amsthm,amssymb,graphicx,enumitem,listofitems}
% Include links
\usepackage[hidelinks]{hyperref}
\usepackage{lmodern}

%%%%%%%%%%%%%  THEOREMS  %%%%%%%%%%%%%%%%
\theoremstyle{plain}
\newtheorem{theorem}{Theorem}
\newtheorem{lemma}[theorem]{Lemma}
\theoremstyle{definition}
\newtheorem*{definition}{Definition}

\theoremstyle{remark}
\newtheorem{remark}{Remark}

% Problem command
\newcounter{problem}
\newcommand{\problem}[2]{%
    \refstepcounter{problem}%
    \noindent\textbf{\href{#1}{Problem~\theproblem.}}#2%
    \vspace{0.3cm}\par%
}

% Solution command
\newcounter{solution}
\newcommand{\solution}[1]{%
    \refstepcounter{solution}%
    \noindent\textbf{Solution.~}#1%
    \vspace{0.3cm}\par%
}

% Solution command that is used when there are more than one solutions.
\newcounter{numsolution}
\newcommand{\numsolution}[2]{%
    \refstepcounter{numsolution}%
    \noindent\textbf{Solution~#1.~}#2%
    \vspace{0.3cm}\par%
}

% Problem command that is used when there are more than one problems with manual numbering.
\newcommand{\numproblem}[2]{%
    \noindent\textbf{Problem~#1.~}#2%
    \vspace{0.3cm}\par%
}

\renewcommand{\mod}{\text{mod }}
\newcommand\cycle[2][\,]{%
\readlist\thecycle{#2}%
(\foreachitem\i\in\thecycle{\ifnum\icnt=1\else#1\fi\i})%
}

%%%%%%%%%%%%%%  PAGE SETUP %%%%%%%%%%%%%%%%
\usepackage[top=1in,left=1in,right=1in,bottom=1.5in]{geometry}
\usepackage{fancyhdr}
\pagestyle{fancy}
\lhead{Generating Functions}
\rhead{\thepage}
\cfoot{}

%%%%%%%%%%%%% SHORTCUTS %%%%%%%%%%%%%%%%%%%%
\newcommand{\half}{\frac{1}{2}}
\newcommand{\cbrt}[1]{\sqrt[3]{#1}}

\begin{document}
\title{Generating Functions}
\author{Anshul Raj Singh}
\date{}
\maketitle

\section{Problems}

\noindent Without talking too much, I will jump right into the problems and try to solve them, because there is no better way to learn math than by doing it. Here is the first problem:\\

\problem{}{
    Show that the number of ways to write a positive integer $n$ as a sum of distinct positive integers is equal to the number of ways to write $n$ as a sum of odd (not necessarily distinct) positive integers.
}

\noindent Let's try some examples to get a better understanding. Take $n=5$. Writing it as a sum of distinct positive integers, we have:
\[
5 = 5,\quad 1+4,\quad 2+3.
\]
\noindent For the second part (odd parts, repetition allowed), we have:
\[
5 = 5,\quad 3+1+1,\quad 1+1+1+1+1.
\]
\noindent So we have three ways in both cases.
\[
6 = 6,\quad 1+5,\quad 2+4,\quad 1+2+3,
\]
and the odd partitions are:
\[
6 = 5+1,\quad 3+3,\quad 3+1+1+1,\quad 1+1+1+1+1+1.
\]
\noindent Again, the counts match. Seems quite reasonable why someone could have conjectured this. But how does one actually prove it?\footnote{There is an elementary solution, but it is not easy to come up with. What we will do instead will feel more natural once you get used to the technique.}

\bigskip

\noindent Before we try to solve it, let's look at another problem.\\

\problem{}{
    Can there be integers $a_1 < a_2 < \cdots < a_n$ such that the differences $a_i - a_j$, $i > j$, take on all values in 
    \[
    \{1,2,3,\ldots,\tfrac{n(n-1)}{2}\}?
    \]
}

\noindent This one looks more complicated, so let's just take a small value and play around.

\noindent Take $n=4$. Try to find such numbers yourself before reading on.

\noindent After some trial and error, you will find that $0,1,4,6$ works. The differences are:
\[
1, 4, 6, 3, 5, 2,
\]
so we get all numbers from $1$ to $6$.

\noindent So we understand the problem now, but again, how do we approach something like this in general?

\bigskip

\noindent If you think about both problems for a while, you will notice the same issue--we do not have a good way to represent these sums and differences and keep track of them at the same time. We need a better way to \emph{encode} this information.

\section{Generating Functions}
Let me get the definition out of the way first.
\begin{definition}
An infinite power series $a_0 + a_1x + a_2x^2 + \cdots$ is called a \textbf{generating series} for the sequence $a_0, a_1, a_2, \ldots$. When we can write this series in a compact form, we call it a \textbf{generating function} for the sequence.
\end{definition}

\begin{remark}
Two power series are equal if and only if the coefficients of $x^n$ are equal for all $n$. This is why generating functions are useful: they allow us to turn counting problems into algebraic identities.
\end{remark}

\noindent This is not something out of the ordinary; I'm sure you have seen this before. For example, the generating function for the sequence $1, 1, 1, \ldots$ is $(1-x)^{-1}$, because
\[
\frac{1}{1-x} = 1 + x + x^2 + \cdots.
\]
Here is another example: the generating function for the sequence $0, 1, 0, 1, 0, 1, \ldots$ is $x(1-x^2)^{-1}$, because
\[
\frac{x}{1-x^2} = x + x^3 + x^5 + \cdots.
\]

\bigskip

\noindent At this point, there is one small technical issue that we should at least mention.

\noindent We have been writing infinite sums like
\[
1 + x + x^2 + x^3 + \cdots = \frac{1}{1-x},
\]
and treating them like ordinary algebraic objects. But strictly speaking, this only makes sense when the series actually \emph{converges}. As you can clearly see, if you put $x=2$, the series becomes $1 + 2 + 4 + 8 + \cdots$, which diverges to infinity but the right side is $-1$. So we need to be careful about when we can actually use these two sides interchangeably.

\noindent Roughly speaking, convergence means that if we keep adding more and more terms, the sum approaches a well-defined value. For example, the series above converges when $|x| < 1$.

\bigskip

\noindent For our purposes, we will simply assume $|x| < 1$ and not worry too much about these details. Under this assumption, all the manipulations we perform with these series and products are justified.

\bigskip

\noindent In fact, in many combinatorial problems, we treat these series purely formally, without worrying about convergence at all.

\bigskip

\noindent Here is another example. Consider the product
\[
(*)\quad A(x) = (1+x)(1+x^2)(1+x^3)\cdots.
\]

\noindent Each bracket $(1+x^k)$ gives us a choice:
\begin{itemize}
\item choose $1$ (ignore the number)
\item choose $x^k$ (include the number $k$).
\end{itemize}

\noindent So when we expand this product, every term we get corresponds to choosing some distinct numbers and adding them up.

\bigskip

\noindent In particular, the coefficient of $x^n$ tells us how many ways we can form $n$. To see this clearly, take $n=5$. The term $x^5$ can appear in the following three ways:
\begin{itemize}
    \item $x^5 = x^5 \longleftrightarrow 5 = 5$,
    \item $x \cdot x^4 = x^{1+4} \longleftrightarrow 1+4 = 5$,
    \item $x^2 \cdot x^3 = x^{2+3} \longleftrightarrow 2+3 = 5$.
\end{itemize}
These correspond exactly to
\[
5,\quad 1+4,\quad 2+3.
\]

\noindent And the coefficient of $x^5$ in $A(x)$ is $3$, which matches the number of ways to write $5$ as a sum of distinct positive integers. So the coefficient of $x^n$ in $A(x)$ counts the number of ways to write $n$ as a sum of distinct positive integers.

\begin{remark}
    Each bracket in $(*)$ is a generating function in itself where all but two of the coefficients are zero. Because we have already established that we can do the standard manipulations with these series, we can multiply them together and get a new generating function. We won't have any issues as we already have taken $|x| < 1$.
\end{remark}


\section{Solutions}

We now have a function $A(x)$ in which the coefficient of $x^n$ counts the number of ways to write $n$ as a sum of distinct positive integers. Now we do the same thing for odd numbers. Try finding this generating function yourself!\\

\noindent Here is the generating function for odd numbers
\[
B(x) = (1 + x + x^{2} + x^{3} + \cdots)(1 + x^3 + x^{6} + x^{9} + \cdots)\cdots = \prod_{k=1}^{\infty} \frac{1}{1-x^{2k-1}}.
\]

\noindent The coefficient of $x^n$ in this expansion counts the number of ways to write $n$ as a sum of odd positive integers (with repetition allowed). Convince yourself that this is true for $n = 10$. How many times is the sum $1+9$ counted?

\bigskip

\noindent So now we have reduced the problem to showing that (why?)
\[
A(x) = B(x).
\]

\noindent We compute
\begin{align*}
\frac{A(x)}{B(x)} &= [(1+x)(1+x^2)(1+x^3)\cdots]\cdot[(1-x)(1-x^3)(1-x^5)\cdots].
\end{align*}

\noindent Notice that $B(x)$ only contains factors of the form $(1 - x^{\text{odd}})$, so we begin by pairing all the odd terms. Note that
\[
(1+x^k)(1-x^k) = 1 - x^{2k}.
\]

\noindent So we pair factors as follows:
\begin{itemize}
    \item $(1+x)$ with $(1-x)$ gives $1-x^2$
    \item $(1+x^3)$ with $(1-x^3)$ gives $1-x^6$
    \item $(1+x^5)$ with $(1-x^5)$ gives $1-x^{10}$
\end{itemize}

\noindent This takes care of all the odd powers. The even powers $(1+x^2),(1+x^4),(1+x^6),\dots$ are left over. So we get
\[
(1-x^2)(1-x^6)(1-x^{10})\cdots \cdot (1+x^2)(1+x^4)(1+x^6)\cdots.
\]

\noindent Now we repeat the same idea again, pair the terms with same powers.

\bigskip

\noindent If we keep doing this, we see that everything keeps cancelling in this way, pushing the powers higher and higher. As this process continues, every factor eventually gets paired and absorbed into higher powers. Since we are working with formal power series (or with $|x|<1$), these higher powers become negligible in the limit, leaving only $1$.

\bigskip

\noindent So
\[
\frac{A(x)}{B(x)} = 1,
\]
and hence $A(x)=B(x)$. This shows that the coefficients of $x^n$ in $A(x)$ and $B(x)$ are equal, which is exactly what we wanted to prove. $\square$

\bigskip

\noindent I will write the solution formally just to show how short it becomes once everything is set up.\\

\numsolution{1}{
Let $A(x) = (1+x)(1+x^2)(1+x^3)\cdots$ and $B(x) = \prod_{k=1}^{\infty} \frac{1}{1-x^{2k-1}}$. The coefficient of $x^n$ in $A(x)$ counts partitions into distinct parts, and the coefficient of $x^n$ in $B(x)$ counts partitions into odd parts. Since 
\begin{align*}
\frac{A(x)}{B(x)} &= [(1+x)(1+x^2)(1+x^3)\cdots]\cdot[(1-x)(1-x^3)(1-x^5)\cdots]\\
&= (1-x^2)(1-x^6)(1-x^{10})\cdots \cdot (1+x^2)(1+x^4)(1+x^6)\cdots\\
&\hspace{0.2cm}\vdots\\
&= 1,
\end{align*}
$A(x)=B(x)$ showing the coefficients are equal. $\square$
}
\vspace{5cm}

\noindent Here are a few exercises to get used to this idea.\footnote{Write some code to verify your answers.}

\begin{enumerate}[label=(\alph*)]
    \item Find the number of positive integer solutions of $x+y+z=10$ using generating functions. What if $x,y,z$ are allowed to be zero?
    \item Find the number of solutions of $x+y+z=50$ such that $x$ is even, $y$ is odd and $z$ is a multiple of $3$.
\end{enumerate}

\section{The Second Problem (A Glimpse)}

\noindent Now let's briefly come back to the second problem. The idea is to build a generating function
\[
A(x) = \sum_{i=1}^{n} x^{a_i}
\]

\noindent and then look at
\[
A(x)A(x^{-1}),
\]
which encodes all differences $a_i - a_j$.

With some clever manipulation (and a bit of complex numbers), one can show that such a configuration cannot exist for $n \geq 5$.

\bigskip

\noindent I will admit that this solution is quite technical and not very intuitive at first glance. Still, I will include it at the end of this document for completeness. I would encourage you to read through it and see how much you can understand, but do not worry if you cannot follow everything. Do not let it intimidate you--these kinds of arguments take time and experience to fully appreciate.

\section{Solutions to Problem 2}
Our first job is to find a generating function that represents the differences $a_i - a_j$, $i > j$. Try it yourself before reading on. The simplest function we know is $A(x) = \sum_{i=1}^{n} x^{a_i}$. After a bit of thought, you will see that, if you multiply $A(x)$ by $A(x^{-1})$, you will get a generating function that has the differences as the exponents of $x$. We have:
\begin{align*}
    A(z)A(z^{-1}) &= \sum_{\substack{i, j = 1\\i > j}}^{n} z^{a_i - a_j} + n + \sum_{\substack{i, j = 1\\i < j}}^{n} z^{a_i - a_j}.
\end{align*}
I have replaced $x$ with $z$, as we'll now be working with complex numbers. Now, look at the first sum: if the problem statement is true, then this sum will be $z^1 + z^2 + \cdots + z^N$, where $N = \frac{n(n-1)}{2} = \binom{n}{2}$, because there are exactly $\binom{n}{2}$ pairs $(i, j)$ with $i > j$. Also, to simplify the expression, note that the second sum is just the first sum with $z$ replaced by $z^{-1}$. So, we have
\begin{align*}
    A(z)A(z^{-1}) &= z^1 + z^2 + \cdots + z^N + n + z^{-1} + z^{-2} + \cdots + z^{-N}\\
    &= \sum_{k = -N}^{N} z^k + n - 1\\
    &= \frac{z^{N+1} - z^{-N}}{z - 1} + n - 1. \quad (\text{Summing the geometric series.})
\end{align*}
To simplify, let us take $z$ to be on the unit circle; that is, $z = e^{i\theta}$ for some $\theta$. The left side becomes $A(z)A(z^{-1}) = A(z)\overline{A(z)} = |A(z)|^2$. Whereas, on the right side, after multiplying the numerator and the denominator by $z^{-1/2}$, we have
\begin{align*}
    \frac{z^{N+1} - z^{-N}}{z - 1} + n - 1 = \frac{z^{N+\frac{1}{2}} - z^{-N-\frac{1}{2}}}{z^{\frac{1}{2}} - z^{-\frac{1}{2}}} + n - 1 = \frac{\sin((N+\frac{1}{2})\theta)}{\sin(\frac{\theta}{2})} + n - 1.
\end{align*}
We used the formula $e^{ix} - e^{-ix} = 2i\sin(x)$ to get the last equality. Therefore, we have
\begin{align*}
    |A(z)|^2 = \frac{\sin((N+\frac{1}{2})\theta)}{\sin(\frac{\theta}{2})} + n - 1 = \frac{\sin(\frac{n^2-n+1}{2}\theta)}{\sin(\frac{\theta}{2})} + n - 1.
\end{align*}

\noindent This equation seems a bit complicated. If we were to find such integers $a_1 < a_2 < \cdots < a_n$ that satisfy the problem statement, then there is no way this equation would help us find them, because we've already gotten rid of all the information about the $a_i$'s. But if these integers do exist for some $n$, then, going back to the starting point, we will have $\binom{n}{2}$ equations and $n$ variables, i.e., these $\binom{n}{2}$ equations, differences take on consecutive values, which seems like it shouldn't be true for large $n$, even though it worked for $n = 4$.

\noindent Although surprising stuff happens all the time, we should at least check if this equation can even be true for all $\theta$ when $n$ is large.\\

\noindent Indeed, it turns out that this equation does not hold for all $\theta$ and large $n$. One simple check you can do is just graph it. You'll see that this starts to fail for $n \geq 5$. But we can also show this analytically. Take $\theta = \frac{3\pi}{n^2-n+1}$. In that case:
\begin{align*}
    \sin\left(\frac{\theta}{2}\right) < \frac{\theta}{2} &\Rightarrow \frac{1}{\sin(\theta/2)} > \frac{2}{\theta}\\
    &\Rightarrow -\frac{1}{\sin(\theta/2)} < -\frac{2}{\theta} = -\frac{(2n^2-2n+2)}{3\pi}.
\end{align*}
We want to see when $-(2n^2-2n+1)/3\pi < -(n-1)$, which will give us the required contradiction. Well, if you check, you'll find that this happens for all $n \geq 5$. So, for $n \geq 5$, there do not exist integers $a_1 < a_2 < \cdots < a_n$ such that the differences $a_i - a_j$, $i > j$, take on all values in the set $\{1, 2, 3, \ldots, \frac{n(n-1)}{2}\}$. $\square$

\bigskip

\section{Practice Problems}

\numproblem{1}{
    Look up how to solve recurrence relations using generating functions.
}

\numproblem{2}{
    Find the number of non negative integer solutions of $x_1 + x_2 + \cdots + x_k = n$ using generating functions.
}

\numproblem{3}{
    Use generating functions to show that every positive integer can be written in exactly one way as a sum of distinct powers of $2$.
}

\numproblem{4}{
    Find a closed form for the generating function
    \[
        \sum_{n=0}^{\infty} \binom{n+k}{k} x^n.
    \](\textbf{Don't differentiate!})
}

\numproblem{5}{
    Let $A=\{a_1, \dots, a_n\}$ and $B=\{b_1, \dots, b_n\}$ be two subsets of positive integers with $A\ne B$ but $A+A=B+B$ where $A+A=\{a_i+a_j \mid a_i, a_j\in A\}$. The number of repetitions in each set is also the same. Show that $n$ is a perfect power of $2$.
}
% \section{References}

\begin{thebibliography}{99}

\bibitem{Newman}
Newman, D. J. Analytic Number Theory. Springer.

\end{thebibliography}

\end{document}